\input{./._relpath-to-latexroot.ltx}
\documentclass[\latexroot/\projectname]{subfiles}
\input{\latexroot/@resources/subfile-setup.ltx}
\begin{document}

\section{Conclusion}\whenintegrated{\label{conclusion}}
\whenintegrated{\label{sec:conclusion}} % integrated doc: SST for labels

Most of the sovereign-debt literature focuses on the \emph{level} of government debt: how much a country can sustain, and at what cost. This paper has shifted the question to the \emph{allocation} of that debt across holder types. The data on advanced economies show that countries with nearly identical debt-to-GDP ratios place their debt in very different hands---some lean on domestic pension funds and insurers, others on foreign investors, and others on their banking sector---and that the bank share rises sharply once total debt passes a threshold of roughly 100--120\% of GDP. Our paper provides a simple, structural account of this heterogeneous pecking order.

On the \emph{model} side, we studied a cost-minimizing government that allocates debt across three sectors, each with a distinct supply structure: domestic non-banks, whose capacity is tied to a country-specific structural sector size $S_h$ and whose cost rises in the utilization rate $d_h/S_h$; foreigners, whose capacity is bounded by a uniform no-default ceiling $\bar{D}_f$ at which the cost of additional borrowing diverges; and banks, which serve as the residual absorber and face a convex welfare cost reflecting crowding-out and financial-repression distortions. The resulting equilibrium is a pecking order whose ordering across sectors depends on country-specific parameters rather than being universal.

On the \emph{empirical} side, taking the model to panel data on advanced economies delivers four findings that map one-to-one to the model's predictions. First, cross-country heterogeneity in debt allocation is driven primarily by the structural size of the non-bank sector, not by the total debt level or by openness alone. Second, the banking sector is the elastic margin: banks absorb the residual after non-bank and foreign capacities are exhausted, which is why the bank share accelerates above the threshold. Third, foreign absorption exhibits the concavity consistent with a uniform no-default ceiling---foreign holdings rise with debt but at a decreasing rate, with a common capacity limit across countries. Fourth, financial integration has only a modest additional effect: more open economies place somewhat less domestic government debt with non-banks, but the effect is second-order relative to the structural sector-size channel.

These results carry a simple message for the policy debate on sustainable debt levels. Two countries with the same debt-to-GDP ratio can face very different sustainability pressures depending on the structural size of their non-bank financial sector and their distance to the foreign no-default ceiling. Countries with large, deep non-bank sectors can absorb more debt at home without triggering the convex repression cost on banks, while countries with structurally small non-bank sectors hit that margin sooner. Our framework does not attempt to micro-found the non-bank demand curve in the~\citet{AiyagariMcGrattan1998} tradition, nor to endogenize sovereign default in the full~\citet{Broneretal2010} or~\citet{GennaioliMartinRossi2014} sense; these extensions, together with richer dynamics of how $\bar{D}_f$ responds to shocks, are natural next steps. Even so, the tractable three-sector structure developed here provides a transparent lens through which to interpret cross-country and over-time variation in who holds government debt, and a portable empirical design for future work on sovereign-debt composition.

% Smart bibliography: Only include bibliography if standalone AND has citations
\smartbib

\pdfonly{
% execute this version if compiling with pdflatex
\end{document}
}
